
Our work sits at the intersection of two research areas: data augmentation techniques for deep learning and learning under label noise. While both areas are well-established, neither addresses the core question of semantic validity---whether standard augmentations violate domain-specific constraints and thereby introduce systematic label noise.

\subsubsection{Data Augmentation Surveys}

Data augmentation has been extensively studied as a regularization technique to improve model generalization. Comprehensive surveys catalogue a wide range of augmentation methods across different modalities: geometric transformations (flip, rotation, crop, affine), photometric adjustments (brightness, contrast, color jitter), learned augmentations (AutoAugment, RandAugment), and mixing strategies (Mixup, CutMix). These surveys typically focus on cataloguing techniques, reporting empirical performance gains, and analyzing computational costs.

However, survey analyses emphasize what augmentations exist and how they improve aggregate performance metrics, not whether augmentations preserve semantic validity for specific classes. The implicit assumption is that standard transformations like horizontal flip are universally applicable---beneficial or at worst neutral---across datasets.

Domain-specific augmentation studies acknowledge that augmentation choices should match domain characteristics. For instance, vertical flip is inappropriate for chest X-rays due to anatomical orientation constraints. Yet these insights remain implicit practitioner knowledge rather than formalized validation frameworks. No prior work systematically tests whether standard augmentations violate semantic constraints on canonical benchmarks like MNIST, or quantifies the resulting degradation.

Automated augmentation search methods discover effective augmentation policies through reinforcement learning or population-based training. However, the search is purely empirical---the learned policies lack explicit semantic validity constraints. If horizontal flip degrades accuracy on asymmetric classes but improves aggregate performance, automated methods may still select it.

\subsubsection{Label Noise Literature}

Learning under label noise addresses training with corrupted ground-truth labels. The standard model distinguishes symmetric noise (labels flipped uniformly across classes) from asymmetric noise (class-specific corruption patterns). Noise-robust training methods include loss correction, sample reweighting, and robust loss functions.

This literature overwhelmingly focuses on annotation errors: mistakes by human labelers, adversarial label flips, or systematic biases in crowdsourced labels. Remarkably, the literature does not consider data augmentation as a source of label noise. The implicit assumption is that labels are corrupted during annotation, not during training through augmentation.

Yet augmentation-induced label noise has distinct properties. First, it is deterministic and tunable: flip probability directly controls the proportion of semantically invalid training examples. Second, it is class-specific: only asymmetric classes are affected by horizontal flip, creating an asymmetric noise pattern. Third, it is reproducible: every training run with flip probability p produces the same noise structure. These properties make augmentation-induced noise an ideal controlled experimental paradigm for testing label noise theories.

Our contribution connects these literatures. We demonstrate that horizontal flip functions as a tunable asymmetric label noise source, with noise rate proportional to flip probability. The perfect or near-perfect dose-response relationships we observe align precisely with label noise theory: accuracy degrades monotonically with noise rate for affected classes.

\subsubsection{MNIST Benchmarks and Practitioner Folklore}

MNIST handwritten digit classification remains a canonical benchmark for validating deep learning techniques. Winning Kaggle solutions on MNIST consistently exclude horizontal flip from their augmentation pipelines, instead favoring rotation ($\pm 10-15$°), small translations, and elastic deformations. PyTorch official tutorials similarly omit horizontal flip while including rotation and affine transformations.

This folklore is well-established but unexplained. Why is flip avoided? Online discussions offer intuition: ''flipping 2 or 5 makes them backwards,'' ''MNIST digits aren't symmetric.'' Yet no published work quantifies the harm, identifies affected classes systematically, or tests the mechanism rigorously.

\subsubsection{Positioning Our Work}

We position our work at the intersection of augmentation surveys and label noise research, filling a gap neither addresses. Augmentation surveys catalogue techniques but lack semantic validity frameworks. Label noise research studies annotation errors but ignores augmentation-induced noise. We demonstrate that:

\begin{enumerate}
\item \textbf{Semantic validity is testable}: Per-class accuracy grouped by digit symmetry reveals differential effects hidden in aggregate metrics.
\end{enumerate}

\begin{enumerate}
\item \textbf{Augmentation induces label noise}: Horizontal flip on MNIST creates class-specific asymmetric noise with dose-response behavior ($\rho$ = -1.0 or -0.969).
\end{enumerate}

\begin{enumerate}
\item \textbf{Folklore has quantitative foundation}: Practitioners correctly avoid flip (0.37-4.10\% degradation on asymmetric digits across flip probabilities {0.3, 0.5, 0.9}).
\end{enumerate}

By demonstrating semantic validity testing on MNIST, we establish a validation methodology applicable beyond MNIST: dose-response testing, differential class-group analysis, and positive control experiments.

